\section{Related Work}\label{sec:related}

\textbf{Poisoning self-updating detectors.} Rubinstein et al.\ showed that a network anomaly detector
that retrains on its own traffic can be walked toward an attack by slowly increasing chaff, the
boiling-frog attack, and defended it with robust PCA on a server~\cite{antidote}. Kloft and Laskov
analysed exactly our detector, the online centroid, and showed that the attacker's share of the training
traffic bounds how far it can be moved; they also proposed switching learning off when the alarm rate
rises~\cite{kloft}. Poisoning of test-time adaptation~\cite{su2025tta}, of on-device learning on
microcontroller-class hardware~\cite{sensors2024}, and of rapid-response adaptation
pipelines~\cite{rapidpoison} shows the problem is current. A 2026 survey of about seventy on-device
learning works for TinyML found no security mechanism among them~\cite{tinymlsurvey}. FOA targets this
gap with a per-device detector that adapts on the device.

\textbf{Adaptive intrusion detection and its gates.} Kitsune made online, per-packet anomaly
detection practical but stops learning after training~\cite{kitsune}. Later adaptive detectors decide
\emph{when} to learn: OWAD detects normality shift and adapts after a test~\cite{owad}; DriftXpert
gates adaptation on an outlier ratio~\cite{driftxpert}; RTTAD updates only from confidently normal
points and stops otherwise~\cite{rttad}; a continual-learning analysis of poisoning proposes rolling back
when the score spikes~\cite{clpoison}; ACORN-IDS and CND-IDS continue learning from normal-only
streams~\cite{acorn,cndids}. Each has a condition under which it stops learning. Our freeze experiments
show that an attacker who can make such a condition true with loud traffic can turn honest drift into a
lasting flood of false alarms. Korycki and Krawczyk discussed attacks that impair adaptation in data
streams~\cite{korycki}; we found no prior evaluation of this freeze effect on the gating rules of
adaptive intrusion detectors.

\textbf{Limits of a single stream.} Recent theory (all 2026 preprints) shows that one label-free score
stream cannot separate honest drift from slow contamination (WARDEN~\cite{warden}), that regularised
continual learners cannot converge under frequent shifted poison~\cite{clpoison}, and that a single
statistic forces a choice between accepting injected changes and rejecting genuine ones
(CAPTURE~\cite{capture}). A further preprint frames self-calibration as an evidence budget with blind
sets~\cite{selfcal}; we have read only its abstract. FOA does not contradict these results; it changes
the observation model by adding peers, and Proposition~\ref{prop:frontier} states what remains impossible.

\textbf{Collaborative detection and robust aggregation.} The closest prior work is
CIoTA~\cite{ciota}: devices of the same model train models collaboratively and accept an update only
when the group agrees, discarding the rest. FOA differs in three ways: peers are chosen to be
vulnerability-\emph{diverse}, so a monoculture botnet cannot manufacture consensus; the remainder is
\emph{rationed} through a budget rather than discarded, which keeps a device able to follow its own
changes; and the guarantee is a bound on movement and time to compromise, not agreement. Byzantine-robust
aggregation (Krum~\cite{krum}, the geometric median~\cite{lopuhaa,minsker}) protects a shared model in
federated learning; FOA uses the geometric median only to estimate the common component of per-device
updates, and keeps every device's own detector.
